% Round24 writer-facing short insert for the Figure 5 boundary chain.
% Nature Communications-safe wording only; no claim beyond the locked evidence chain.

\paragraph{Figure 5 boundary summary.}
A scalar sidecar bridge is repeatable, but it does not survive the Figure-5-like shared-carrier protocol; backend decoupling improves ordinary separation but does not restore the non-common margin, and control-family refinement does not recover stable common-specificity under the current criterion.

\begin{table}[t]
\centering
\small
\setlength{\tabcolsep}{4pt}
\begin{tabular}{p{0.10\linewidth}p{0.18\linewidth}p{0.18\linewidth}p{0.42\linewidth}}
\hline
Step & CMO & CMN/control & Locked role \\
\hline
Round14 & $0.786 \pm 0.050$ & $0.788 \pm 0.031$ & Repeatable scalar sidecar bridge \\
Round15 & $0.312 \pm 0.026$ & $0.348 \pm 0.043$ & Shared-carrier failure boundary \\
Round16 & $0.542 \pm 0.024$ & $0.378 \pm 0.057$ & Backend decoupling improves ordinary only \\
Round17 & $0.542 \pm 0.024$ & $0.464 \pm 0.082$ & No stable common-specific recovery under current criterion \\
\hline
\end{tabular}
\caption{Extreme-short quantitative anchor table for the Figure 5 boundary chain. Here CMO denotes common-minus-ordinary and CMN/control denotes common-minus-non-common control.}
\end{table}

\paragraph{SI short summary.}
Across Round14--17, the audited boundary chain narrows from a repeatable scalar sidecar bridge to a shared-carrier failure boundary, then shows that backend readout decoupling and control-family refinement do not rescue a stable non-common margin under the current criterion.

\paragraph{Reviewer-safe robustness note.}
The boundary chain remains directionally stable under the Round19 audit: $0/4$ alternate anchors changed the manuscript role, whereas the Round17 control-family step remained fully stable for only $3/9$ audited threshold pairs. For that reason, the Round17 sentence should remain ``does not recover stable common-specificity under the current criterion'' rather than an absolute impossibility claim.
